\chapter{2020:Emmanuelle Charpentier and Jennifer A. Doudna}

\label{chap:chemistry2020}

The Nobel Prize in Chemistry for the year 2020 was awarded jointly to Emmanuelle Charpentier and Jennifer A. Doudna.

\textbf{Motivation:}

for the development of a method for genome editing

\section{Emmanuelle Charpentier}

\subsection*{Early Life and Education}

Emmanuelle Charpentier was born in 1968 in Juvisy-sur-Orge, France.

\subsection*{Career and Major Research}

Charpentier received her doctorate from the Institut Pasteur in Paris in 1995. She held research positions in Europe and the United States and became scientific director at the Max Planck Unit for the Science of Pathogens in Berlin. Her work on the bacterial immune system CRISPR showed that it could be programmed by RNA to cut DNA at specific sites.

\subsection*{Nobel Prize-winning Work}

The work for which Emmanuelle Charpentier was awarded the Nobel Prize in Chemistry in 2020 was described by the Nobel Committee as: "for the development of a method for genome editing".

In 2012 Charpentier and Jennifer Doudna, with Martin Jinek, showed that the protein Cas9, guided by a dual RNA, could cut any DNA sequence chosen by the researcher.

\subsection*{Significance and Impact}

This reprogrammable RNA-guided nuclease became the basis of the CRISPR-Cas9 method, which has made genome editing fast, simple, and available to laboratories worldwide.

\subsection*{Later Career and Honors}

Charpentier continues to direct research on RNA biology and bacterial infection at her institute in Berlin.

\subsection*{Legacy}

CRISPR-Cas9 genome editing is used in fundamental research, agriculture, and the development of therapies for genetic disease.

\subsection*{Selected Publications}

"A Programmable Dual-RNA-Guided DNA Endonuclease in Adaptive Bacterial Immunity" (Science, 2012)

\clearpage

\section{Jennifer A. Doudna}

\subsection*{Early Life and Education}

Jennifer A. Doudna was born in 1964 in Washington, D.C., USA.

\subsection*{Career and Major Research}

Doudna earned a bachelor's degree from Pomona College in 1985 and a doctorate from Harvard University in 1989. She became a professor at the University of California, Berkeley, and an investigator of the Howard Hughes Medical Institute, working on the structure and function of RNA.

\subsection*{Nobel Prize-winning Work}

The work for which Jennifer A. Doudna was awarded the Nobel Prize in Chemistry in 2020 was described by the Nobel Committee as: "for the development of a method for genome editing".

With Emmanuelle Charpentier, Doudna established how the CRISPR-associated protein Cas9 is programmed by guide RNA and demonstrated its programmable DNA cleavage in 2012.

\subsection*{Significance and Impact}

Her biochemical characterization explained how Cas9 recognizes and cuts its target, providing the understanding needed to turn the system into a general genome-editing tool.

\subsection*{Later Career and Honors}

Doudna continues her research at the University of California, Berkeley, and is a co-founder of companies that apply CRISPR technology.

\subsection*{Legacy}

CRISPR-Cas9 has transformed genetic research and is being developed as a treatment for inherited diseases and other conditions.

\subsection*{Selected Publications}

"A Programmable Dual-RNA-Guided DNA Endonuclease in Adaptive Bacterial Immunity" (Science, 2012)

\clearpage

\section*{Chapter Summary}

The 2020 prize recognized the development of CRISPR-Cas9, a method for genome editing that originated in a bacterial immune system. Charpentier and Doudna showed that a single RNA-guided protein can cut any chosen DNA sequence, and the technique is now used across biology and medicine.
